\clearpage\fancyhead[L]{\small GAUSS AN BODE.}\hypertarget{native-362-lower}{}\label{native-362-lower}
\begin{center}\textbf{GAUSS AN BODE.}\par\vspace{4pt}\hrule\vspace{4pt}\fontsize{10}{12}\selectfont Astronomisches Jahrbuch für 1815. Seite 245\;..\;249. Berlin 1812.\par\vspace{4pt}\hrule\vspace{4pt}\end{center}

\begin{flushright}Göttingen den 22. August 1812.\end{flushright}
\par\indent Die erste Auffindung der Pallas habe ich Ihnen, wenn ich nicht irre, bereits
angezeigt. Die Anzahl der hiesigen Beobachtungen ist diesmal nur klein. Um
die Zeit der Opposition war der Stand des
\SourceLeaf{363}Planeten für unser Locale zu hoch, um ihn noch mit dem gewöhnlichen
Instrument am Kreismicrometer zu observiren; ich habe mich daher eines
andern, schlechtern Instruments bedienen müssen, und auch so hat mir
schlechtes Wetter und Mondschein nur Eine Beobachtung verstattet. Meine
sämmtlichen Beobachtungen sind folgende:

\begin{center}\fontsize{10}{12}\selectfont\setlength{\tabcolsep}{4pt}\renewcommand{\arraystretch}{1.15}\begin{tabular}{l|c|c|c}
& Mittlere Zeit & Scheinb. Aufsteig. & Abweichung N.\\\hline
April 9 & $11^{\mathrm u}\;32^{\mathrm m}\;27^{\mathrm s}$ & $269^\circ\;5'\;58''4$ & $16^\circ\;38'$\\
Mai 2 & $10\;36\;6$ & $268\;42\;5.3$ & $21\;4\;31''7$\\
— 4 & $10\;44\;44$ & $268\;31\;38.0$ & $21\;25\;4.3$\\
Juni 7 & $11\;13\;51$ & $262\;50\;51.1$ & $25\;2\;56.6$\\
\end{tabular}\end{center}

\par\indent Von auswärtigen Beobachtungen sind mir ausser zwei frühern vom Hrn. Dr.
\textsc{Olbers}, nur vier Meridianbeobachtungen des Herrn Dr. \textsc{Burckhardt} auf der
Sternwarte der Militairschule zu Paris mitgetheilt, die ich in Nr. 127
unsrer gel. Anz. [S. 356 d. B.] habe abdrucken lassen. Aus diesen Datis hat
Hr. \textsc{Nicolai}, dessen Geschicklichkeit im Rechnen Sie bereits kennen, die
Opposition berechnet:

\Needspace{10\baselineskip}1812 Juni 10. $3^{\mathrm u}\;2^{\mathrm m}\;48^{\mathrm s}$ Mittlere Zeit in Göttingen

\begin{center}\fontsize{10}{12}\selectfont\setlength{\tabcolsep}{2pt}\renewcommand{\arraystretch}{1.15}\begin{tabular}{lcr}
Wahre Länge & $.\;.\;.\;.\;.\;.\;.\;.\;.\;.\;.$ & $259^\circ\;27'\;51''7$\\
Wahre geoc. Breite & $.\;.\;.\;.\;.\;.\;.\;.\;.$ & $48\;16\;26.1\;\mathrm{N}.$\\
\end{tabular}\end{center}

\Needspace{13\baselineskip}\par\indent Nachdem nur kleine Correctionen an die letzten Elemente angebracht waren,
nemlich

\begin{center}\fontsize{10}{12}\selectfont\setlength{\tabcolsep}{4pt}\renewcommand{\arraystretch}{1.15}\begin{tabular}{rl}
$-13''29$ & bei der Epoche der Länge 1810\\
$-0.01200$ & bei der tägl. Bewegung\\
$-17.57$ & bei der Länge des Perihels\\
$-4.09$ & beim Excentricitätswinkel\\
$+0''80$ & bei der Länge des Knoten\\
$-1.21$ & bei der Neigung der Bahn\\
$+0.0000045$ & beim Log. der halben grossen Axe;\\
\end{tabular}\end{center}

\Needspace{17\baselineskip}\noindent war die möglich genaueste Uebereinstimmung derselben folgende.

\begin{center}\fontsize{10}{12}\selectfont\setlength{\tabcolsep}{6pt}\renewcommand{\arraystretch}{1.15}\begin{tabular}{c|c|c}
& \multicolumn{2}{c}{Differenz}\\
Oppos. & Mittl. Länge & hel. Breite\\\hline
$1803$ & $-5''06$ & $-2''78$\\
$1804$ & $-12.95$ & $-9.66$\\
$1805$ & $-15.20$ & $-0.67$\\
$1807$ & $+8.28$ & $-5.77$\\
$1808$ & $+16.04$ & $+11.67$\\
$1809$ & $+13.00$ & $-1.48$\\
$1811$ & $+11.73$ & $-9.53$\\
$1812$ & $-25.92$ & $+6.15$\\
\end{tabular}\end{center}

\par\indent Ich schreibe diese Unterschiede hauptsächlich den Störungen durch Saturn
und Mars zu, da bisher blos die Störungen durch den Jupiter in Betracht
gezogen wurden. Letztere sind indess bei weitem die erheblichsten. Die
erste Berechnung derselben hat mir ungefähr 400 Gleichungen gegeben. Bei
einer zweiten vollständigern Rechnung, wozu ich einen kleinen Anfang
gemacht habe, werden ungefähr 1000 entwickelt werden, die über $0''1$
betragen. Die nächste künftige Opposition wird nach Hrn. \textsc{Nicolai}'s
Rechnung eintreten

\Needspace{10\baselineskip}1813 August 18. $8^{\mathrm u}\;50^{\mathrm m}\;5^{\mathrm s}$ Mittlere Zeit in Göttingen

\begin{center}\fontsize{10}{12}\selectfont\setlength{\tabcolsep}{2pt}\renewcommand{\arraystretch}{1.15}\begin{tabular}{lcr}
Wahre Länge & $.\;.\;.\;.\;.\;.\;.\;.\;.\;.\;.\;.$ & $325^\circ\;24'\;18''$\\
Wahre geocentrische Breite & $.\;.\;.\;.\;.\;.\;.$ & $24\;37\;24\;\mathrm{N}..$\\
\end{tabular}\end{center}

\Needspace{17\baselineskip}\SourceLeaf{364}\par\indent Die Beobachtungen der \emph{Juno} waren in diesem Jahre überaus schwierig,
und schon das blosse Auffinden war in der von kleinen Sternen wimmelnden
Milchstrasse bei der grossen Lichtschwäche des Planeten, der nur etwa die
11te, und später noch nicht die 10te Grösse erreichte, mit vieler Mühe
verbunden. Ich habe indess doch das Glück gehabt, nachdem ich sie einmal
herausgefunden hatte, 4 gute Beobachtungen zu machen, die ich Ihnen hier
mittheile.

\begin{center}\fontsize{10}{12}\selectfont\setlength{\tabcolsep}{4pt}\renewcommand{\arraystretch}{1.15}\begin{tabular}{l|c|c|c}
1812 & Mittlere Zeit & Scheinb. Aufsteig. & Abweichung\\\hline
Jan. 7 & $12^{\mathrm u}\;21^{\mathrm m}\;0^{\mathrm s}$ & $282^\circ\;21'\;24''5$ & $4^\circ\;41'\;44''6\;\mathrm S.$\\
13 & $12\;51\;12$ & $281\;16\;32.9$ & $4\;33\;39.6$\\
Juli 12 & $11\;35\;28$ & $275\;5\;10.3$ & $5\;6\;11.2$\\
13 & $11\;28\;12$ & $274\;52\;49.9$ & $5\;9\;20.0$\\
\end{tabular}\end{center}

\par\indent Nach einer mir eigenthümlichen Methode, die entferntere Beobachtungen mit
anzuwenden erlaubt, als das gewöhnliche Verfahren, und deren Erklärung
ich mir auf eine andere Gelegenheit aufsparen muss, hat aus diesen
Beobachtungen Hr. \textsc{Wachter} die Opposition folgendermassen bestimmt:

\Needspace{10\baselineskip}1812 Juni 29. $19^{\mathrm u}\;51^{\mathrm m}\;43^{\mathrm s}$ Mittlere Zeit in Göttingen

\begin{center}\fontsize{10}{12}\selectfont\setlength{\tabcolsep}{2pt}\renewcommand{\arraystretch}{1.15}\begin{tabular}{lcr}
Wahre Länge & $.\;.\;.\;.\;.\;.\;.\;.\;.\;.\;.\;.$ & $278^\circ\;15'\;21''3$\\
Wahre geocentrische Breite & $.\;.\;.\;.\;.\;.\;.$ & $18\;36\;56.9$\\
\end{tabular}\end{center}

\Needspace{14\baselineskip}\par\indent Die Unterschiede meiner letzten Elemente von den Beobachtungen sind auf
gewöhnliche Weise angestellt. Es sind nach Hrn. \textsc{Wachter}'s Berechnung
folgende:

\begin{center}\fontsize{10}{12}\selectfont\setlength{\tabcolsep}{7pt}\renewcommand{\arraystretch}{1.15}\begin{tabular}{l|c|c}
& AR. & Declin.\\\hline
Juni 7 & $+21''8$ & $-35''9$\\
13 & $+31.0$ & $-34.7$\\
Juli 12 & $+22.4$ & $-46.9$\\
13 & $+26.2$ & $-45\;0$\\
\end{tabular}\end{center}

\Needspace{16\baselineskip}\par\indent Derselbe geschickte Rechner hat ferner die Elemente der 4 letzten
Oppositionen nach meiner im I. Bande der Neuen Comment. der hiesigen
K. Societät [S. 1. d. B.] erklärten Methode angepasst, und folgende
Resultate gefunden:

\begin{center}\fontsize{9}{11}\selectfont\setlength{\tabcolsep}{2pt}\renewcommand{\arraystretch}{1.15}\begin{tabular}{lcr}
Epoche, Merid. von Göttingen & $.\;.\;.\;.\;.\;.\;.\;.$ & $177^\circ\;48'\;21''0$\\
Tägliche mittlere tropische Bewegung & $.\;.\;.\;.\;.$ & $813''25748$\\
Länge des Perihels 1811 & $.\;.\;.\;.\;.\;.\;.\;.\;.\;.\;.$ & $53^\circ\;15'\;10''1$\\
Länge des $\AscendingNode$ & $.\;.\;.\;.\;.\;.\;.\;.\;.\;.\;.\;.\;.\;.$ & $171\;9\;16.7$\\
Neigung der Bahn & $.\;.\;.\;.\;.\;.\;.\;.\;.\;.\;.\;.\;.$ & $13\;4\;17.2$\\
Logarithm. der halben grossen Axe & $.\;.\;.\;.\;.\;.$ & $0.4265679$\\
\end{tabular}\end{center}

\par\indent Endlich hat er auch auf das nächste Jahr die Ephemeride für die Juno
berechnet, wovon ich das Vergnügen habe, Ihnen eine Copie beizulegen. Es
scheint übrigens nicht, dass die Juno in diesem Jahre noch anderswo
beobachtet wäre.

\par\indent Die Vesta hat Hr. Prof. \textsc{Harding} ein paarmal am Kreismikrometer beobachtet:
die Ephemeride gab die AR. etwa $11'$ zu klein. Ohne Zweifel wird ihre
$\text{\AstroSymbols ☍}$
sehr gut durch Meridianbeobachtungen bestimmt werden können. Eher habe ich
keine neue Rechnungen darüber vornehmen lassen wollen; auch wird es
hinreichend sein, wenn eine neue Ephemeride etwa mit dem Schluss des Jahres
1813 anfängt, daher also für das Jahrb. 1815 noch keine nöthig sein
wird.

\SourceLeaf{365}\par\indent Auch die Ceres ist zur Zeit der
$\text{\AstroSymbols ☍}$
hier einigemal in und ausser dem Meridian beobachtet; es sind indess noch
keine Rechnungen darüber angestellt.

\par\indent Den in den Zeitungen angekündigten neuen Cometen haben Prof. \textsc{Harding} und
ich einen Abend umsonst gesucht.


